\documentclass[11pt]{article}
\usepackage{amsmath,amsthm,amssymb}
\usepackage[margin=1in]{geometry}
\usepackage{booktabs,hyperref}

\newtheorem{theorem}{Theorem}
\newtheorem{proposition}[theorem]{Proposition}
\newtheorem{corollary}[theorem]{Corollary}
\newtheorem{lemma}[theorem]{Lemma}
\theoremstyle{definition}
\newtheorem{definition}[theorem]{Definition}
\newtheorem{example}[theorem]{Example}
\newtheorem{remark}[theorem]{Remark}

\newcommand{\Hom}{\operatorname{Hom}}
\newcommand{\Ext}{\operatorname{Ext}}
\newcommand{\Aut}{\operatorname{Aut}}
\newcommand{\id}{\mathrm{id}}
\newcommand{\ob}{\operatorname{ob}}
\newcommand{\Ab}{\mathbf{Ab}}
\newcommand{\Sk}{\mathrm{Sk}}
\newcommand{\C}{\mathcal{C}}
\newcommand{\D}{\mathcal{D}}
\newcommand{\Z}{\mathbb{Z}}
\newcommand{\CatB}{\mathbf{Cat}_B}
\newcommand{\Igp}{\mathbf{I}}
\newcommand{\Bpi}{B\pi}
\newcommand{\infl}{\mathrm{inf}}
\newcommand{\res}{\mathrm{res}}
\newcommand{\tg}{\mathrm{tg}}
\newcommand{\im}{\operatorname{im}}

\title{One spectral sequence, two obstructions:\\
the Zappa--Sz\'ep holonomy is the transgression, not a summand}
\author{MacBeth\\\small Kodamai}
\date{17 September 2026}

\begin{document}
\maketitle

\begin{abstract}
The $H^2$-cluster carries two proved degree-two obstruction theories on a
connected groupoid: the multiplicative Zappa--Sz\'ep \emph{composition holonomy}
$[\omega]\in H^2(\Sk_\C;\D)$ over the orbit base, and the additive
direction-functor class $H^2_{\CatB}(G;A)\cong H^2(\pi;M)$ of the vertex group,
the home of Ferri prunability. A natural conjecture (the ``product
decomposition'') hoped these were two independent summands of one categorical
cohomology, $H^2_{\mathrm{cat}}(G;M)\cong H^2(\pi;M)\oplus H^2_{\mathrm{base}}(\Sk_\C;M)$,
with both possibly nonzero. \textbf{We prove that this clean direct sum is
impossible precisely when it would be interesting, and we replace it by the
correct structure.} When the Zappa--Sz\'ep factorisation is a normal one
$1\to\D\to K\to\Gamma\to1$ (loop-bundle fibre $\D$, orbit base
$\Gamma=\Sk_\C$), the categorical cohomology is the abutment of the
Lyndon--Hochschild--Serre spectral sequence $E_2^{p,q}=H^p(\Gamma;H^q(\D;M))$.
Its two edges are exactly the two cluster faces --- the base edge
$E_2^{2,0}=H^2(\Gamma;M^\D)$ carrying $[\omega]$ and the fibre edge
$E_2^{0,2}=H^2(\D;M)^\Gamma$ carrying prunability --- and the transgression
$d_2\colon H^1(\D;M)^\Gamma\to H^2(\Gamma;M^\D)$ is cup with the holonomy,
$d_2(\varphi)=\varphi_*[\omega]$. Consequently: (i) the clean $\oplus$ holds iff
the sequence collapses at $E_2$ with vanishing cross term --- which forces
$[\omega]=0$ (split extension) --- and in the connected/coarse-base case
$\Gamma\simeq\ast$ it recovers $H^2_{\mathrm{cat}}=H^2(\pi)$ exactly (Prop.~A,
T1); (ii) \textbf{no-go}: if $[\omega]\neq0$ and some $\varphi\colon\D\to M$ has
$\varphi_*[\omega]\neq0$, the base edge does not survive and no direct sum
exists. Thus the two obstructions are never simultaneously present as
independent summands: a nonzero base class \emph{is} the differential that
transgresses itself away. This turns the three separate cluster results into
\emph{one} object --- a spectral sequence whose transgression is the ZS
holonomy --- refining the proved separation $[\omega]\perp$ prunability from
``orthogonal'' to ``the two edges coupled by $[\omega]$''. All finite claims are
machine-checked ($\Z/4$: $\dim H^2=1$ vs.\ naive $3$; $Q_8$: $2$ vs.\ $6$;
$\Z/2\times\Z/2$: $3=3$; $\Z/6$: $1=1$).
\end{abstract}

\section{The two faces and the conjecture they invite}

Three proved results place two obstruction theories on one connected groupoid
$G$ with vertex group $\pi=G(x_0,x_0)$.

\begin{itemize}
\item[\textbf{(H)}] \emph{Holonomy face.} For the multiplicative groupoid of a
quiver skew brace factored (Zappa--Sz\'ep) over a vertex bundle $\D$, the
global-closure obstruction is a degree-two class $[\omega]\in H^2(\Sk_\C;\D)$
over the orbit base $\Sk_\C$; the factorisation $K=\C\bowtie\D$ exists iff
$[\omega]=0$ \cite{gobstr}. When $\D$ is normal this is the Schreier extension
class of $1\to\D\to K\to\Gamma\to1$, $\Gamma=\Sk_\C$, and it is genuinely
nonzero, e.g.\ $[\omega]\neq0$ for $\Z/4$ over $\Z/2$ \cite{gzs}.
\item[\textbf{(P)}] \emph{Prunability face.} The direction-functor cohomology of
Ambra--Duvieusart--Montoli \cite{adm} reduces, for a connected groupoid, to the
vertex-group cohomology, $H^2_{\CatB}(G;A)\cong H^2(\pi;M)$ (T1 of \cite{sep2});
this is the additive face that carries Ferri's prunability obstruction
\cite{sep}.
\item[\textbf{(S)}] \emph{Separation.} $[\omega]\perp$ prunability: the two are
orthogonal, proved from a computed witness \cite{sep} and re-derived
structurally from the decomposition $G\cong\Igp(\ob G)\times\Bpi$ (T2 of
\cite{sep2}).
\end{itemize}

The tempting synthesis, recorded as the target \texttt{h2cat-product-decomposition},
is a Künneth-style direct sum making both faces visible at once:
\begin{equation}\label{eq:conj}
   H^2_{\mathrm{cat}}(G;M)\;\overset{?}{\cong}\;
   \underbrace{H^2(\pi;M)}_{\text{prunability}}\;\oplus\;
   \underbrace{H^2_{\mathrm{base}}(\Sk_\C;M)}_{[\omega]},
   \qquad\text{both summands possibly nonzero.}
\end{equation}

\paragraph{The conflation that breaks it.} Statement \eqref{eq:conj} silently
uses one symbol $\Sk_\C$ for two different bases.
\begin{itemize}
\item In the \emph{connected/full-loop-bundle} regime of T2, $\Sk_\C=\Igp(\ob G)$
is the \emph{coarse groupoid on the objects}. A coarse groupoid is a disjoint
union of codiscrete (contractible) components, so $H^{\ge1}(\Sk_\C;-)=0$
\emph{always}: the base summand of \eqref{eq:conj} is identically zero
(Prop.~A). Here $G\cong\Igp(\ob G)\times\Bpi$ is a genuine categorical product,
Künneth applies, and it says nothing but $H^2_{\mathrm{cat}}=H^2(\pi)$.
\item In the \emph{proper-normal-bundle} regime of \cite{gzs}, the base
$\Sk_\C=B\Gamma=B(K/\D)$ is the classifying groupoid of a \emph{quotient group}
and is generally \emph{not} contractible; this is where $[\omega]\neq0$. But now
$G$ is \emph{not} a product $\Igp\times B\Gamma$: the data
$(\D,\Gamma,\pi)$ form a non-split \emph{extension} $1\to\D\to\pi\to\Gamma\to1$.
\end{itemize}
So the two regimes never coexist inside one product. The honest object that
sees a non-contractible base is not a categorical product but a group
\emph{extension}, and its cohomology is governed by the
Lyndon--Hochschild--Serre (LHS) spectral sequence, not by Künneth. We prove that
in this correct object $[\omega]$ appears as the \emph{transgression
differential}, which is exactly why \eqref{eq:conj} fails whenever the base
class is nonzero.

\section{Setup and the reduction to a group extension}

We work in the normal Zappa--Sz\'ep regime of \cite{gzs}, the only regime in
which the base carries a nonzero class.

\begin{proposition}[Reduction, \cite{gzs} Prop.~2.1--2.2, \cite{sep2} T1]
\label{prop:reduce}
Let $G$ be a connected groupoid whose vertex group $\pi$ admits a normal abelian
subgroup $\D\trianglelefteq\pi$ (the vertex bundle of a Zappa--Sz\'ep
factorisation), with quotient $\Gamma=\pi/\D$ the vertex group of the orbit base
$\Sk_\C=B\Gamma$. Then:
\begin{enumerate}
\item[\textup{(a)}] $\D$ is a $\Gamma$-module via conjugation, and the
Zappa--Sz\'ep holonomy $[\omega]\in H^2(\Sk_\C;\D)=H^2(\Gamma;\D)$ is the
Schreier class of the extension
\begin{equation}\label{eq:ext}
   1\longrightarrow\D\longrightarrow\pi\xrightarrow{\ q\ }\Gamma\longrightarrow1;
\end{equation}
\item[\textup{(b)}] for a $G$-module $A$ with $M:=A(x_0)$, the categorical
(direction-functor) cohomology reduces to the vertex group,
$H^2_{\mathrm{cat}}(G;A)=H^2_{\CatB}(G;A)\cong H^2(\pi;M)$.
\end{enumerate}
\end{proposition}
\begin{proof}
(a) is \cite[Prop.~2.1--2.2]{gzs}: normality gives $\Sk_\C=\pi/\D=\Gamma$ and the
defect $2$-cochain is the Schreier factor set of \eqref{eq:ext}. (b) is T1 of
\cite{sep2}, $H^n_{\CatB}(G,A)\cong H^n(\pi,M)$, whose $n=2$ instance rests on
the reconstructed Schreier correspondence there and elementary homological
algebra.
\end{proof}

Proposition~\ref{prop:reduce} moves the whole question into ordinary group
cohomology: \emph{how does $H^2(\pi;M)$ decompose along the extension
\eqref{eq:ext}, and where do the two faces $[\omega]$ and $H^2(\D)$ sit?} The
answer is the LHS spectral sequence.

\section{The spectral sequence and its two edges}

\begin{theorem}[LHS decomposition of the categorical cohomology]\label{thm:lhs}
With the data of Proposition~\ref{prop:reduce}, $H^\bullet_{\mathrm{cat}}(G;A)\cong
H^\bullet(\pi;M)$ is the abutment of the first-quadrant cohomological spectral
sequence
\begin{equation}\label{eq:E2}
   E_2^{p,q}\;=\;H^p\bigl(\Gamma;\,H^q(\D;M)\bigr)\;\Longrightarrow\;H^{p+q}(\pi;M),
\end{equation}
of the extension \eqref{eq:ext} \cite[\S6.8]{weibel}. Its two degree-two edge
terms are exactly the two $H^2$-cluster faces:
\[
   E_2^{2,0}=H^2\bigl(\Gamma;M^\D\bigr)
   \quad\text{(\emph{base edge}: the home of $[\omega]$)},
   \qquad
   E_2^{0,2}=H^2(\D;M)^\Gamma
   \quad\text{(\emph{fibre edge}: the prunability face)} .
\]
The abutment is filtered with associated graded
$\mathrm{gr}\,H^2(\pi;M)=E_\infty^{2,0}\oplus E_\infty^{1,1}\oplus E_\infty^{0,2}$,
where $E_\infty^{p,q}$ is the subquotient of $E_2^{p,q}$ surviving the
differentials.
\end{theorem}
\begin{proof}
The LHS spectral sequence of a group extension with abelian coefficient module
$M$ is standard \cite[Thm.~6.8.2]{weibel}; \eqref{eq:E2} is its $E_2$-page. The
edge identifications $E_2^{p,0}=H^p(\Gamma;M^\D)$ and
$E_2^{0,q}=H^q(\D;M)^\Gamma$ are the definition of the two axes. That the base
edge is where $[\omega]$ lives is Proposition~\ref{prop:reduce}(a)
($[\omega]\in H^2(\Gamma;\D)$, mapped to $H^2(\Gamma;M^\D)$ by any $\Gamma$-map
$\D\to M^\D$); that the fibre edge is the vertex/prunability face is
\cite{sep,sep2}. The three-step filtration in total degree $2$ is the standard
convergence of a first-quadrant sequence.
\end{proof}

The point of Theorem~\ref{thm:lhs} is that the two faces are not
\emph{a priori} summands: they are two edges of one page, and whether they
survive to the abutment is decided by the differentials. The relevant one is the
transgression, and it is the holonomy itself.

\begin{lemma}[The transgression is cup with the holonomy]\label{lem:tg}
In the five-term exact sequence of \eqref{eq:E2},
\[
   0\to H^1(\Gamma;M^\D)\xrightarrow{\infl}H^1(\pi;M)\xrightarrow{\res}
   H^1(\D;M)^\Gamma\xrightarrow{\ \tg\ }H^2(\Gamma;M^\D)\xrightarrow{\infl}H^2(\pi;M),
\]
the transgression $\tg=d_2^{0,1}$ is given by pushforward of the extension class:
for a $\Gamma$-equivariant homomorphism $\varphi\in\Hom_\Gamma(\D,M)=H^1(\D;M)^\Gamma$,
\[
   \tg(\varphi)\;=\;\varphi_*[\omega]\;\in\;H^2(\Gamma;M^\D),
\]
where $\varphi_*\colon H^2(\Gamma;\D)\to H^2(\Gamma;M^\D)$ is induced by $\varphi$
and $[\omega]$ is the Schreier class of \eqref{eq:ext}.
\end{lemma}
\begin{proof}
We give the cocycle computation for the case $M$ a trivial $\pi$-module (which
covers all witnesses below and the $\Z/4$-Schreier line); the general
$\Gamma$-module case is \cite[proof of the transgression formula]{nsw}, identical
after keeping the covariant end. Choose a normalised set section
$s\colon\Gamma\to\pi$ of $q$, $s(1)=1$, with Schreier factor set
\[
   s(a)\,s(b)=\omega(a,b)\,s(ab),\qquad \omega(a,b)\in\D,\quad [\omega]=[\omega]\in H^2(\Gamma;\D).
\]
An element of $E_2^{0,1}=H^1(\D;M)^\Gamma$ is a $\Gamma$-invariant homomorphism
$\varphi\colon\D\to M$. Running the construction of the transgression in the
double complex of the extension \cite[\S6.8]{weibel}: $\varphi$ lifts to the
$1$-cochain $\tilde\varphi\in C^1(\pi;M)$, $\tilde\varphi(g)=\varphi(\text{$\D$-part of }g)$
along the section, and its image under the vertical differential, pushed to the
base, is the $2$-cochain
\[
   (a,b)\longmapsto \varphi\bigl(\omega(a,b)\bigr)=(\varphi_*\omega)(a,b).
\]
Thus $\tg(\varphi)=[\varphi_*\omega]=\varphi_*[\omega]$. (Concretely: $\tg(\varphi)$
is the obstruction to extending $\varphi$ from $\D$ to a homomorphism
$\pi\to M$; it is $0$ iff the pushed-out extension
$1\to M\to\varphi_*\pi\to\Gamma\to1$ splits, i.e.\ iff $\varphi_*[\omega]=0$.)
\end{proof}

\section{Collapse: recovering the conjecture where it is true}

\begin{corollary}[When the clean sum holds]\label{cor:collapse}
The naive direct sum \eqref{eq:conj} of the two edges equals $H^2(\pi;M)$
whenever the sequence \eqref{eq:E2} collapses at $E_2$ in total degree $2$ with
vanishing cross term $E_2^{1,1}$. Each of the following is sufficient:
\begin{enumerate}
\item[\textup{(i)}] \textbf{Coarse/connected base} $\Gamma=1$
\textup{(}Prop.~A, T2 of \cite{sep2}\textup{)}: then $E_2^{p,q}=0$ for $p>0$, the
base edge and cross term vanish, and $H^2(\pi;M)=E_2^{0,2}=H^2(\D;M)$. Since here
$\D=\pi$ this is the identity $H^2_{\mathrm{cat}}(G)=H^2(\pi)$ of T1 --- the clean
sum degenerates to its fibre summand, and $[\omega]=0$.
\item[\textup{(ii)}] \textbf{Direct product} $\pi=\D\times\Gamma$: then $[\omega]=0$
and the sequence collapses (Künneth), giving the \emph{full} associated graded
$H^2(\pi;M)\cong E_2^{2,0}\oplus E_2^{1,1}\oplus E_2^{0,2}$ --- the honest
Künneth form, including the cross term $H^1(\Gamma)\otimes H^1(\D)$ that
\eqref{eq:conj} omits.
\item[\textup{(iii)}] \textbf{Coprime orders} $\gcd(|\D|,|\Gamma|)=1$
\textup{(}Schur--Zassenhaus\textup{)}: $E_2^{p,q}=0$ for $p,q>0$, so
$H^2(\pi;M)=E_2^{2,0}\oplus E_2^{0,2}$ and $[\omega]=0$.
\end{enumerate}
In every collapsing case the base class is forced to zero: $[\omega]=0$.
\end{corollary}
\begin{proof}
(i) A coarse groupoid is equivalent to the terminal category, so
$H^{\ge1}(\Gamma;-)=0$; the surviving column is $p=0$. (ii) A direct product
splits \eqref{eq:ext}, so $[\omega]=0$; collapse is the Eilenberg--Zilber/Künneth
theorem for $B\D\times B\Gamma$. (iii) For coprime orders the higher cohomology
of each factor is annihilated by both $|\D|$ and $|\Gamma|$, hence $0$ off the
edges; a split extension (Schur--Zassenhaus) has $[\omega]=0$.
\end{proof}

Corollary~\ref{cor:collapse} is the precise content of the original conjecture:
\eqref{eq:conj} is \emph{true exactly on the locus $[\omega]=0$}, and there it is
either the trivial statement $H^2_{\mathrm{cat}}=H^2(\pi)$ (connected case) or
the genuine Künneth sum (direct product) --- but in neither case is the base
summand nonzero. The interesting ``both nonzero'' regime is empty, as we now
show.

\section{The no-go: a nonzero base class transgresses itself away}

\begin{theorem}[{No clean sum when $[\omega]\neq0$}]\label{thm:nogo}
Suppose $[\omega]\neq0$ and there is a $\Gamma$-equivariant $\varphi\colon\D\to M$
with $\varphi_*[\omega]\neq0$ in $H^2(\Gamma;M^\D)$ \textup{(}e.g.\ $M=\D$,
$\varphi=\id$, whenever $[\omega]\neq0$\textup{)}. Then:
\begin{enumerate}
\item[\textup{(a)}] the transgression $d_2^{0,1}=\tg$ is nonzero
\textup{(}Lemma~\ref{lem:tg}\textup{)}, so the base edge does not survive:
$E_\infty^{2,0}\subsetneq E_2^{2,0}$, and the image $\varphi_*[\omega]$ of the
holonomy is \emph{killed} in the abutment;
\item[\textup{(b)}] consequently $H^2(\pi;M)$ is \emph{not} the direct sum
\eqref{eq:conj} of the two edges: $\dim H^2(\pi;M)<\dim E_2^{2,0}+\dim E_2^{1,1}
+\dim E_2^{0,2}$ whenever $\tg\neq0$.
\end{enumerate}
In particular the two $H^2$-cluster faces are \emph{never} simultaneously present
as independent summands of $H^2_{\mathrm{cat}}(G)$: a nonzero base class is not a
summand but a differential.
\end{theorem}
\begin{proof}
(a) By Lemma~\ref{lem:tg}, $\tg(\varphi)=\varphi_*[\omega]\neq0$, so
$d_2^{0,1}\neq0$. The differential $d_2^{0,1}\colon E_2^{0,1}\to E_2^{2,0}$ has
nonzero image, hence $E_3^{2,0}=E_2^{2,0}/\im d_2^{0,1}$ is a proper quotient;
as the base edge receives no further differentials
($d_r^{p,0}$ into $(2,0)$ come from $(2-r,r-1)$, which for $r\ge3$ is a negative
column), $E_\infty^{2,0}=E_3^{2,0}\subsetneq E_2^{2,0}$, and the class
$\varphi_*[\omega]$ (a generator of $\im d_2$) is annihilated. (b) The abutment
has $\dim H^2=\sum_{p+q=2}\dim E_\infty^{p,q}\le\sum_{p+q=2}\dim E_2^{p,q}$ with
strict inequality once any $E_\infty^{p,q}\subsetneq E_2^{p,q}$; part (a) gives
this at $(2,0)$.
\end{proof}

\begin{corollary}[Witnesses, machine-checked]\label{cor:witness}
For trivial $\Z/2$-coefficients $M=\mathbb F_2$
\textup{(}\texttt{2026-09-17-lhs-decomposition-check.py}\textup{)}:
\begin{center}\small
\begin{tabular}{lcccccl}
\toprule
$\pi=\D.\Gamma$ & split? & $[\omega]$ & $\dim H^2(\pi)$ &
$\dim(E_2^{2,0}\!\oplus\!E_2^{1,1}\!\oplus\!E_2^{0,2})$ & $\mathrm{rk}\,\tg$ & verdict\\
\midrule
$\Z/4=\Z/2.\Z/2$          & no  & $\neq0$ & $1$ & $1{+}1{+}1=3$ & $1$ & \textbf{no sum} (base edge dies)\\
$Q_8=(\Z/2)^2.\Z/2$       & no  & $\neq0$ & $2$ & $3{+}2{+}1=6$ & $1$ & \textbf{no sum}\\
$\Z/2\times\Z/2=\Z/2.\Z/2$& yes & $0$     & $3$ & $1{+}1{+}1=3$ & $0$ & clean (full Künneth)\\
$\Z/6=\Z/2.\Z/3$          & yes & $0$     & $1$ & $0{+}0{+}1=1$ & $0$ & clean (fibre only)\\
\bottomrule
\end{tabular}
\end{center}
For $\Z/4$ the transgression is an isomorphism $E_2^{0,1}\xrightarrow{\sim}E_2^{2,0}$,
so the entire base edge $H^2(\Z/2;\mathbb F_2)$ --- the home of the nonzero Schreier
$[\omega]$ of $\Z/4$ over $\Z/2$ \cite{gzs} --- \emph{vanishes} from the
abutment. The naive sum overcounts by $2$; the true $H^2(\Z/4;\mathbb F_2)$ is
$1$-dimensional.
\end{corollary}

\section{What the cluster is, correctly}

The three cluster results do unify into one object --- but that object is a
spectral sequence, not a direct sum.

\begin{theorem}[Structural unification]\label{thm:unify}
Under Proposition~\ref{prop:reduce}, the categorical cohomology
$H^2_{\mathrm{cat}}(G;A)\cong H^2(\pi;M)$ is the total degree-two term of the LHS
spectral sequence \eqref{eq:E2}, in which:
\begin{enumerate}
\item the \emph{fibre edge} $E_2^{0,2}=H^2(\D;M)^\Gamma$ is the
additive/prunability face \textup{(P)};
\item the \emph{base edge} $E_2^{2,0}=H^2(\Gamma;M^\D)$ is the holonomy face,
receiving the transgression $d_2$;
\item the \emph{transgression} $d_2\colon E_2^{0,1}\to E_2^{2,0}$ \emph{is} the
Zappa--Sz\'ep holonomy, $d_2(\varphi)=\varphi_*[\omega]$ \textup{(H)}.
\end{enumerate}
The proved separation $[\omega]\perp$ prunability \textup{(S)} is the special
case $\Gamma\simeq\ast$ \textup{(}base contractible, only the fibre edge
survives\textup{)}; in general the two faces are not orthogonal summands but the
two axes of one page, \emph{coupled} by the transgression, which is $[\omega]$
itself. The clean direct sum \eqref{eq:conj} holds iff this coupling vanishes,
i.e.\ iff $[\omega]=0$.
\end{theorem}
\begin{proof}
Assemble Theorem~\ref{thm:lhs} (edges), Lemma~\ref{lem:tg} (transgression $=$
holonomy), Corollary~\ref{cor:collapse} (collapse $\Leftrightarrow$ the true
locus of \eqref{eq:conj}) and Theorem~\ref{thm:nogo} (no-go off that locus). The
separation is Corollary~\ref{cor:collapse}(i): $\Gamma=1\Rightarrow$ base edge
and cross term vanish, $H^2_{\mathrm{cat}}=H^2(\D)=H^2(\pi)$, orthogonality is
the triviality of the base axis.
\end{proof}

\begin{remark}[Why this is the right correction]
The conjecture \eqref{eq:conj} failed for a structural reason worth stating
plainly: a nonzero base class $[\omega]$ \emph{is} a non-split extension, and a
non-split extension is exactly a nonzero transgression. One cannot ``add'' the
holonomy as a summand next to the prunability class, because the holonomy is the
mechanism by which the base cohomology is glued into the fibre cohomology.
The separation theorem saw the two faces as orthogonal only because it lived on
the contractible-base slice where the transgression is forced to zero; off that
slice, orthogonality is replaced by a single coupling differential. This is a
strictly more informative statement, and it is the honest ``internal
replacement / Zappa--Sz\'ep'' deliverable: compositional correctness of the
directed-container composite $= [\omega]=0 =$ the collapse of one spectral
sequence, whose two edges are the base handoff geometry and the vertex/additive
prunability.
\end{remark}

\section{Status and honest scope}
\paragraph{Proved.} Theorem~\ref{thm:lhs} (LHS decomposition; standard SS +
proved edge identifications), Lemma~\ref{lem:tg} (transgression $=$ holonomy;
self-contained cocycle proof for trivial coefficients, cited \cite{nsw} in
general), Corollary~\ref{cor:collapse} (collapse locus $=\{[\omega]=0\}$),
Theorem~\ref{thm:nogo} (no-go), Theorem~\ref{thm:unify} (unification). These rest
on the proved Proposition~\ref{prop:reduce}(a)(b) $=$ \cite{gzs,sep2}, the
classical LHS machinery \cite{weibel,nsw}, and Schur--Zassenhaus.
\paragraph{Computed.} The four witnesses of Corollary~\ref{cor:witness}
(machine-checked bar complex over $\mathbb F_2$ and five-term rank count).
\paragraph{Not claimed.} The clean direct sum \eqref{eq:conj} with both summands
nonzero --- \textbf{refuted} (Theorem~\ref{thm:nogo}). The nonabelian /
non-normal Zappa--Sz\'ep matched-pair regime (Kac/Pirashvili $H^2$) is out of
scope, as in \cite{gzs}. The precise \emph{identification} of Ferri prunability
with a specific element of the fibre edge $E_2^{0,2}$ remains the separately
tracked open node.

\begin{thebibliography}{99}
\bibitem{gobstr} MacBeth, \emph{The global-closure obstruction (G) is a
cohomology class}, Kodamai note, 11 June 2026.
\bibitem{gzs} MacBeth, \emph{Invertibility does not kill the Zappa--Sz\'ep merge
obstruction: the groupoid case is group cohomology}, Kodamai note, 24 July 2026
(\texttt{2026-07-24-groupoid-zs-obstruction.tex}).
\bibitem{sep} MacBeth, \emph{Ferri prunability is not the Zappa--Sz\'ep
composition holonomy: a separation theorem}, Kodamai note, 16 September 2026
(\texttt{2026-09-16-ferri-prunability-vs-zs-holonomy.tex}).
\bibitem{sep2} MacBeth, \emph{The direction-functor $H^2$ of a connected
groupoid is its vertex-group cohomology}, Kodamai note, 16 September 2026
(\texttt{2026-09-16-h2catb-vertex-group-reduction.tex}).
\bibitem{adm} S.~Ambra, A.~Duvieusart, A.~Montoli, \emph{A direction functor
approach to the cohomology of small categories}, arXiv:2608.07380 (2026).
\bibitem{weibel} C.~A.~Weibel, \emph{An Introduction to Homological Algebra},
Cambridge Univ.\ Press, 1994 (\S6.8, the Lyndon--Hochschild--Serre spectral
sequence and the five-term exact sequence).
\bibitem{nsw} J.~Neukirch, A.~Schmidt, K.~Wingberg, \emph{Cohomology of Number
Fields}, 2nd ed., Springer, 2008 (Ch.~I, transgression and the extension class).
\bibitem{brown} K.~S.~Brown, \emph{Cohomology of Groups}, GTM 87, Springer, 1982.
\end{thebibliography}
\end{document}
